% Irodov, user-supplied PDF, problems 271--280, printed pp.43--45
% (PDF pp.52--54). Source diagrams 1.71--1.72 and answer sheet checked.

\exproblem{271}{43}
边长为 $a$ 的均匀立方体放在摩擦系数为 $k$ 的
水平面上，获得一初始水平速度后减速至停止。
解释它在滑行减速时关于\emph{固定于水平地面且垂直于运动方向的轴}
的角动量为何逐渐减少并消失，
并求地面支持力合力作用线相对重力作用线的水平偏移
$\Delta x$。
\exsolution
设立方体质量为 $m$，运动方向取为 $+x$，
立方体质心离地高度为 $a/2$。
滑动摩擦力大小为 $f=kmg$，指向 $-x$。
摩擦作用于底面，关于\emph{经过质心}的横向水平轴产生大小
\[
N_f=f(a/2)=kmg(a/2)
\]
的倾覆力矩。若立方体只平动而不倾转，
总转矩须为零。重力作用线穿过质心，
而地面总支持力 $N=mg$ 可以在底面上偏离几何中心，
对该轴产生反向力矩 $N\,\Delta x$。
因此
\[
mg\,\Delta x=kmg\,\frac a2
\quad\Longrightarrow\quad
\boxed{\Delta x=\frac12ka}.
\]
由于地面摩擦使质心减速，立方体关于\emph{固定于地面}
的横向水平轴所具有的平动轨道角动量逐渐减少至零。
支持力作用线偏移只保证关于\emph{质心}的总力矩为零、
不发生自转，并不能阻止平动角动量的变化。
\excheck 压力中心的偏移源自底面压力分布的不对称；
若 $\Delta x>a/2$，即 $k>1$，上述不倾转状态
不再自洽，应考虑倾翻。

\exproblem{272}{43}
一根长为 $\ell$、质量为 $M$ 的均匀水平细杆 $AB$
可绕通过端点 $A$ 的竖直固定轴无摩擦旋转，
最初角速度为 $\omega_0$。
质量为 $m$ 的小套环从轴心 $A$ 开始无摩擦地
沿杆滑向另一端 $B$。
求套环到达 $B$ 时相对杆的径向滑动速率 $v'$。
\exsolution
杆对端点的转动惯量为
$I_0=M\ell^2/3$。
套环离轴距离为 $r$ 时，总角动量
$L_z=(I_0+mr^2)\omega$。
外力对固定竖轴的力矩为零，故
\[
(I_0+mr^2)\omega=I_0\omega_0.
\]
套环在 $r=\ell$ 时
$\omega_f=I_0\omega_0/(I_0+m\ell^2)$。
无摩擦且轴不作功，机械能守恒：
\[
\tfrac12I_0\omega_0^2
=\tfrac12(I_0+m\ell^2)\omega_f^2
+\tfrac12m{v'}^2.
\]
代入 $\omega_f$：
\[
m{v'}^2
=I_0\omega_0^2
\left[1-\frac{I_0}{I_0+m\ell^2}\right]
=\frac{I_0m\ell^2\omega_0^2}{I_0+m\ell^2}.
\]
因此
\[
\boxed{v'=\frac{\omega_0\ell}{\sqrt{1+3m/M}}}.
\]
\excheck $v'$ 是套环沿杆的\emph{相对径向速率}，
不是套环对地总速率，后者还含方位速度分量。

\exproblem{273}{43}
质量为 $m=5.0\,\mathrm{kg}$、
长度为 $\ell=90\,\mathrm{cm}$ 的均匀细杆
在光滑水平面上静止，一端突然受到
大小为 $J=3.0\,\mathrm{N\,s}$、
垂直于杆的水平冲量。
求随后的运动中，杆一半对另一半的内力大小。
\exsolution
冲量使杆的质心获得平动速度
$v_C=J/m$；关于质心的冲量矩为
$J\ell/2$。
均匀细杆对质心的转动惯量为
$I_C=m\ell^2/12$，所以碰后角速度
\[
\omega=\frac{J\ell/2}{I_C}
=\frac{6J}{m\ell}.
\]
冲量结束后外部水平合力为零，
质心匀速运动、角速度恒定。
每半杆的质心距整杆质心 $\ell/4$，
做圆周运动所需向心加速度为
$\omega^2\ell/4$。
半杆质量 $m/2$，沿杆指向整体质心的
内力大小因此为
\[
F_{\rm internal}=\frac m2\cdot\frac{\omega^2\ell}{4}
=\frac{m\ell\omega^2}{8}.
\]
代入得
\[
\boxed{F=\frac{9J^2}{2m\ell}
=\frac{9(3.0)^2}{2(5.0)(0.90)}
=9.0\,\mathrm N}.
\]
\excheck 该内力是在\emph{碰撞结束后的转动阶段}
维持两半杆的向心加速度，
不是瞬时冲量 $J$ 除以未知碰撞时间得到的力。

\exproblem{274}{43}
质量为 $M$、边长为 $\ell$ 的均匀正方形薄板，
绕其一条竖直侧边形成的固定轴无摩擦转动。
质量为 $m$ 的小球以速度 $v$ 垂直薄板撞击其中心，
发生完全弹性碰撞。
求：(a) 小球碰后有符号速度 $v'$；
(b) 碰后薄板轴承受到的水平合力大小。
\exsolution
薄板对过边的竖直轴惯量
$I=\frac13M\ell^2$，
撞击点距轴 $b=\ell/2$。
设碰撞时小球对薄板冲量为 $J$，
则小球碰后速度 $v'=v-J/m$，
薄板角速度
\[
\omega=\frac{Jb}{I}.
\]
完全弹性的一维局部接触恢复条件为
$b\omega-v'=v$，即
$J(b^2/I+1/m)=2v$。
因 $b^2/I=3/(4M)$，
\[
J=\frac{2v}{1/m+3/(4M)}
=\frac{8mMv}{4M+3m}.
\]
于是
\[
\boxed{v'=v\frac{3m-4M}{3m+4M}}.
\]
碰后板质心距转轴 $\ell/2$，
它以角速度 $\omega$ 作圆周运动，
所需向心加速度为 $\omega^2\ell/2$，
因此轴的水平合反力大小为
\[
F=M\omega^2\frac\ell2,\qquad
\omega=\frac{J\ell/2}{M\ell^2/3}
=\frac{12mv}{\ell(4M+3m)}.
\]
故
\[
\boxed{F=\frac{72Mm^2v^2}{\ell(4M+3m)^2}
=\frac{8Mv^2}{\ell[1+4M/(3m)]^2}}.
\]
\excheck $v'$ 可为负，意味着小球反弹；
碰撞时角动量关于轴守恒，但整体线动量会受轴冲量，
不能直接令两物体的线动量守恒。

\exproblem{275}{43--44}
一根质量为 $M$、长度为 $\ell$ 的均匀细杆
上端悬挂在固定水平铰点，初始竖直静止。
质量为 $m$ 的子弹水平飞来击中杆下端并嵌入，
随后杆摆过最大角度 $\alpha$。设 $m\ll M$。
求：(a) 子弹飞来速率；
(b) 碰撞过程中“子弹+杆”系统\emph{线动量增量}
及其原因；(c) 为使碰撞过程中系统线动量不变，
子弹应击中距上端多远的点 $x$？
\exsolution
碰撞很短，关于铰点的冲量矩为零，
故子弹撞前角动量 $mv\ell$ 等于碰后杆子系统的
角动量 $I\omega$，其中
$I\simeq M\ell^2/3$，从而
$\omega\simeq3mv/(M\ell)$。
碰后摆起阶段机械能守恒。
最大角处杆的质心升高
$\frac\ell2(1-\cos\alpha)$，
忽略 $m/M$ 高阶修正，
\[
\frac12I\omega^2
=\frac12Mg\ell(1-\cos\alpha).
\]
联立求得
\[
\boxed{v\simeq\frac M m
\sqrt{\frac{2g\ell}{3}}\sin\frac\alpha2}.
\]
碰后线动量主要来自杆的质心运动：
$P_f\simeq M(\omega\ell/2)
=\frac32mv$，碰前为 $P_i=mv$。
故碰撞过程中的水平线动量增量
\[
\boxed{\Delta p\simeq\frac12mv
=\frac M2\sqrt{\frac{2g\ell}{3}}\sin\frac\alpha2
=M\sqrt{\frac{g\ell}{6}}\sin\frac\alpha2}.
\]
它由固定铰点施加的\emph{外部水平冲量}造成。
若击中距上端为 $x$ 的位置，
碰后 $\omega\simeq mvx/I$，
线动量 $P_f\simeq M(\ell/2)\omega$。
令 $P_f=mv$ 得
\[
\boxed{x=\frac{2I}{M\ell}\simeq\frac{2\ell}{3}}.
\]
\excheck 此处始终采用 $m\ll M$；
撞击阶段不用机械能守恒，摆起阶段才使用。

\exproblem{276}{44}
一根质量为 $M$、半径为 $R$ 的均匀水平圆盘
可绕过中心的固定竖轴自由旋转。
盘上径向导槽中有质量为 $m$ 的小物体，
以穿过轴心的轻绳相连。
初始物体位于盘边缘，整体以角速度 $\omega_0$ 转动。
随后缓慢拉绳使物体移到轴心。
求：(a) 最后角速度；(b) 拉绳外力做功。
\exsolution
盘的转动惯量
$I_d=\frac12MR^2$。
拉绳通过竖轴，不提供关于轴的外力矩，
故总角动量守恒：
\[
(I_d+mR^2)\omega_0=I_d\omega_f.
\]
由此
\[
\boxed{\omega_f=\left(1+\frac{2m}{M}\right)\omega_0}.
\]
缓慢操作使最终径向振动动能可忽略，
拉绳的外力功等于系统转动动能增量：
\[
W=\frac12I_d\omega_f^2
-\frac12(I_d+mR^2)\omega_0^2.
\]
将 $\omega_f$ 代入：
\[
\boxed{W=\frac12mR^2\omega_0^2
\left(1+\frac{2m}{M}\right)}.
\]
\excheck 角动量守恒但动能\emph{增加}，
多出的能量来自拉绳外力做功。

\exproblem{277}{44}
均匀水平圆盘质量为 $m_2$、半径为 $R$，
能绕中心固定竖轴无摩擦转动。
质量为 $m_1$ 的人站在盘边缘。
原先整体静止，此后人沿盘边行走，
相对于圆盘累计转过角度 $\varphi'$，
相对于盘面的切向速率为 $v'(t)$，最后停下。
忽略人的尺寸，求：
(a) 圆盘相对地面转过角度 $\varphi$；
(b) 人对圆盘施加的竖轴力矩 $N(t)$。
\exsolution
圆盘惯量
$I_d=\frac12m_2R^2$。
人在地面系中的角速度为
$\dot\varphi+\dot\varphi'$。
因整个系统关于竖轴无外力矩，初始角动量零：
\[
I_d\dot\varphi+
m_1R^2(\dot\varphi+\dot\varphi')=0.
\]
所以
$\dot\varphi=
-\frac{m_1R^2}{I_d+m_1R^2}\dot\varphi'$。
由零初角积分：
\[
\boxed{\varphi=-\frac{2m_1}{2m_1+m_2}\varphi'}.
\]
人对圆盘的力矩等于盘角动量变化率：
\[
N=I_d\ddot\varphi
=-\frac{I_d m_1R^2}{I_d+m_1R^2}\ddot\varphi'.
\]
因 $v'=R\dot\varphi'$，
\[
\boxed{N(t)=-\frac{m_1m_2R}{2m_1+m_2}
\frac{\mathrm dv'}{\mathrm dt}}.
\]
\excheck 人加速向前走时盘反向转动；
当人减速并最终相对盘停住时两者地面系角速度
也同时回到零，但盘可能已有非零角位移。

\exproblem{278}{44}
两个水平圆盘绕同一固定竖轴自由旋转，
惯量分别为 $I_1,I_2$，初始角速度为
$\omega_1,\omega_2$。
上盘落到下盘，最终因摩擦一起以同一角速度旋转。
求：(a) 最终角速度；(b) 此过程摩擦力做的总功。
\exsolution
轴无外力矩，故总角动量守恒：
\[
I_1\omega_1+I_2\omega_2=(I_1+I_2)\omega.
\]
因此
\[
\boxed{\omega=\frac{I_1\omega_1+I_2\omega_2}{I_1+I_2}}.
\]
两圆盘之间的滑动摩擦使动能损失，
其对两盘的\emph{合计做功}等于末初动能差：
\[
W_f=\frac12(I_1+I_2)\omega^2
-\frac12I_1\omega_1^2-\frac12I_2\omega_2^2.
\]
代入公共角速度并整理：
\[
\boxed{W_f=-\frac{I_1I_2}{2(I_1+I_2)}
(\omega_1-\omega_2)^2}.
\]
\excheck $W_f\le0$，且仅当
$\omega_1=\omega_2$ 时等于零；
角动量守恒不表示转动动能守恒。

\exproblem{279}{44}
一小圆片质量为 $m$，在光滑水平面上以速度 $v$
沿垂直于一根静止均匀细杆的方向运动，
随后与该杆的一端发生完全弹性碰撞。
细杆长为 $\ell$、质量为 $\eta m$，
其中 $\eta>0$。求碰后圆片速度 $v'$、
细杆角速度 $\omega$；
并求圆片碰后停下或反弹的质量比条件。
\exsolution
设沿入射方向的碰撞冲量大小为 $J$。
圆片碰后有符号速度
$v'=v-J/m$。
杆的质心平动速率
$V=J/(\eta m)$。
杆关于质心的惯量
$I=\eta m\ell^2/12$，
碰撞点距质心为 $\ell/2$，
故其角速度大小
\[
\omega=\frac{J(\ell/2)}{I}
=\frac{6J}{\eta m\ell}.
\]
杆端沿碰撞方向的速率为
$V+\omega\ell/2=4J/(\eta m)$。
弹性碰撞恢复系数为1，因而两接触点
碰后的相对分离速率等于碰前接近速率：
\[
\frac{4J}{\eta m}-\left(v-\frac Jm\right)=v.
\]
所以 $J=\frac{2\eta m v}{4+\eta}$，
\[
\boxed{v'=v\frac{4-\eta}{4+\eta}},
\qquad
\boxed{\omega=\frac{12v}{\ell(4+\eta)}}.
\]
\excheck 当 $\eta=4$，$v'=0$；
当 $\eta>4$，$v'<0$，圆片沿反方向弹回；
$0<\eta<4$ 时圆片仍向前运动。

\exproblem{280}{45}
图1.72的水平平台 $P$ 可绕固定竖轴 $OO'$
自由旋转；平台、电动机和配重关于此轴的惯量为 $I$。
电动机装有轻框架，框架支承一颗可沿自身轴 $BB'$
自由自转的均匀球 $A$；球关于任一直径的惯量为 $I_0$。
初始 $BB'$ 与 $OO'$ 重合，
球以自转角速度 $\omega_0$ 旋转，平台原先静止。
电动机慢慢将自转轴 $BB'$ 从竖直方向转过
$90^\circ$ 或 $180^\circ$，求：
(a) 两种转角所需电动机功；
(b) 轴转至 $90^\circ$ 后，为保持平台轴竖直
所需外力矩的大小。
\exsolution
设球自转轴与竖直方向的瞬时夹角为 $\theta$，
平台绕竖直轴的角速度为 $\Omega$。
由于球绕自身轴\emph{自由}自转，
平台自转对球附加的角速度
仅为\emph{垂直于}球自转轴的分量，
其大小为 $\Omega\sin\theta$。
因此关于竖直轴的总角动量为
\[
L_z=[I+I_0\sin^2\theta]\Omega
+I_0\omega_0\cos\theta.
\]
系统竖直轴上无外力矩，初始
$L_z=I_0\omega_0$，所以
\[
\Omega(\theta)=
\frac{I_0\omega_0(1-\cos\theta)}
{I+I_0\sin^2\theta}.
\]
转轴缓慢改变时忽略其瞬态动能，
相对于初始纯自转的附加动能为
\[
\Delta K=\frac12[I+I_0\sin^2\theta]\Omega^2
=\frac{I_0^2\omega_0^2(1-\cos\theta)^2}
{2[I+I_0\sin^2\theta]}.
\]
无耗散时该增量等于电动机做功，故
\[
\boxed{W_{90^\circ}=
\frac{I_0^2\omega_0^2}{2(I+I_0)}},\qquad
\boxed{W_{180^\circ}=\frac{2I_0^2\omega_0^2}{I}}.
\]
在 $\theta=90^\circ$ 时平台角速度为
$\Omega=I_0\omega_0/(I+I_0)$。
此时球的主自转角动量
$I_0\omega_0$ 位于水平面，
该矢量随平台以 $\Omega$ 绕竖直轴转动，
保持轴线方向所需外力矩大小等于其方向变化率：
\[
\boxed{N=I_0\omega_0\Omega
=\frac{I_0^2\omega_0^2}{I+I_0}}.
\]
\excheck 若误认为球与平台绕竖直轴\emph{刚性锁定}，
就会在 $\theta=180^\circ$ 得到错误的转动惯量；
原题明确说球可绕 $BB'$ 自由转动，
必须区分自转轴向和平面正交分量。
